% !PRINTED 305
\weberpage{305}

\begin{center}
\textbf{§ 82. \gesperrt{Divisors of $\mathfrak{L}_0$ of index $p$ for $p = 5, 7, 11$.}}
\end{center}

There remains the possibility that for $p = 5, 7, 11$ divisors of $\mathfrak{L}_0$ of index $p$ exist.

1. $p = 5$. We investigate first the case $p = 5$. Since $m$ and $m - 1$ cannot be congruent to 0, there remain for $m$ the assumptions
\[
m \equiv 2, 3, 4 \quad (\mathrm{mod}\ p).
\]
If, however, we replace the substitution $W$ § 81, (21) by
\weq{1}{WB = \begin{pmatrix} -\xi\varrho, & \xi \\ -\xi^{-1}m, & \xi^{-1}\varrho^{-1}(m-1) \end{pmatrix},}
whereby $m$ passes into $1 - m$, then the case $m \equiv 4$ is brought back to the case $m \equiv 2$.

For $m \equiv 2$, however, (23), § 81 is not satisfied, and of the congruences (22) none is possible, since $\pm 2$ are non-residues of 5 and the left sides of (22), 3., 4. reduce for $m = 2$ to $\xi^4 - 1$, which for every $\xi$ different from zero is divisible by 5. It therefore only remains that
\weq{2}{m \equiv 3, \quad m' \equiv 3 \quad (\mathrm{mod}\ 5)}
and the two systems $W, W'$ [§ 81, (16), (17)] can differ from one another only in this, that $\alpha^{-1}\beta$, and hence also $\varrho$, is in the one a quadratic residue, in the other a quadratic non-residue of 5. The group sought therefore consists, if it exists, of the 12 substitutions contained in the three forms
\weq{3}{\begin{pmatrix} \xi, & 0 \\ 0, & \xi^{-1} \end{pmatrix}, \quad \begin{pmatrix} 0, & \xi \\ -\xi^{-1}, & 0 \end{pmatrix}, \quad \begin{pmatrix} \xi, & \eta \\ 2\eta^{-1}, & 3\xi^{-1} \end{pmatrix}}
in which
\[
\xi \equiv 1, 2, \quad \eta \equiv \pm 1, \pm 2 \quad (\mathrm{mod}\ 5)
\]
is to be set.

2. $p = 7, 11$.

Since $-1$ is to be found among the quadratic non-residues for $p = 7, 11$, the composed substitution $WB$, (1), belongs to the $W'$ [because the quadratic character of the quantity denoted by $\varrho$ in (18), § 81 is the opposite one in $W$ and in $WB$]. Accordingly, by (1) and § 81, (19), (21), $\beta'\gamma' \equiv -m \equiv m' - 1$, hence:
\weq{4}{m + m' \equiv 1 \quad (\mathrm{mod}\ p).}
